\documentclass[11pt,a4paper]{article}
\usepackage[utf8]{inputenc}
\usepackage[T1]{fontenc}
\usepackage[french]{babel}
\usepackage{geometry}
\geometry{margin=2cm}
\usepackage{amsmath,amssymb}
\usepackage{tikz}

\title{Exercice -- Lecture graphique d'images et d'antécédents}
\author{}
\date{}

\begin{document}

\maketitle

On considère la fonction $f$ définie sur l'intervalle $[-4 ; 5]$, dont la courbe représentative $\mathcal{C}_f$ est tracée dans le repère orthogonal ci-dessous. La courbe passe par les points $(-4 ; 2)$, $(-2 ; 4)$, $(0 ; 0)$, $(2 ; -3)$, $(3 ; 0)$ et $(5 ; 4)$.

\begin{center}
\begin{tikzpicture}[scale=0.9]
    \draw[->, >=stealth, gray] (-5,0) -- (6,0) node[right] {$x$};
    \draw[->, >=stealth, gray] (0,-4) -- (0,5) node[above] {$y$};
    \foreach \x in {-4,-3,-2,-1,1,2,3,4,5}
        \draw[gray!50] (\x,-3.5) -- (\x,4.5);
    \foreach \y in {-3,-2,-1,1,2,3,4}
        \draw[gray!50] (-4.5,\y) -- (5.5,\y);
    \draw[thick, blue, domain=-4:5, samples=100] plot (\x, {0.25*(\x)^3 - 1.25*\x});
    \node[blue, above right] at (3,1) {$\mathcal{C}_f$};
    \filldraw[red] (-4,2) circle (2pt) node[above left] {$(-4;2)$};
    \filldraw[red] (-2,4) circle (2pt) node[above left] {$(-2;4)$};
    \filldraw[red] (0,0) circle (2pt) node[below right] {$(0;0)$};
    \filldraw[red] (2,-3) circle (2pt) node[right] {$(2;-3)$};
    \filldraw[red] (3,0) circle (2pt) node[above] {$(3;0)$};
    \filldraw[red] (5,4) circle (2pt) node[above left] {$(5;4)$};
\end{tikzpicture}
\end{center}

\section*{Partie A -- Lectures d'images}
En utilisant exclusivement le graphique représentant la fonction $f$ :
\begin{enumerate}
    \item Déterminer l'image de $0$ par la fonction $f$.
    \item Déterminer l'image de $-2$ par la fonction $f$.
    \item Déterminer l'image de $3$ par la fonction $f$.
    \item Le nombre $4$ a-t-il pour image $5$ ou pour antécédent $5$ sur le graphique ? Justifier graphiquement.
\end{enumerate}

\section*{Partie B -- Lectures d'antécédents}
En utilisant exclusivement le graphique représentant la fonction $f$ :
\begin{enumerate}
    \item Déterminer le ou les antécédents de $0$ par la fonction $f$.
    \item Déterminer le ou les antécédents de $4$ par la fonction $f$.
    \item Le nombre $-3$ possède-t-il des antécédents ? Si oui, les déterminer graphiquement.
    \item Résoudre graphiquement l'équation $f(x) = -3$.
\end{enumerate}

\end{document}
